\input{../tex-inputs/latex-leseansicht-vorspann}

\section[Gustav Schwarzkopf an Arthur Schnitzler, 16. 8. 1901]{L04303 Gustav Schwarzkopf an Arthur Schnitzler, 16. 8. 1901}
\nopagebreak\mylabel{L04303v}
\rehead{ }\normalsize\beginnumbering\briefempfaengerindex{Schnitzler, Arthur@\textsc{Schnitzler, Arthur}!zzzSchwarzkopf, Gustav@\emph{von Gustav Schwarzkopf}!1901-08-161@{16. 8. 1901}|(be}
\toendnotes[C]{\smallbreak\pagebreak[2]}
\correspDesc{Versand  durch Gustav Schwarzkopf am 16. 8. 1901 in Wien
\newline{}Übermittlung  am 16. 8. 1901 in Wien
\newline{}Erhalt  durch Arthur Schnitzler im Zeitraum [17. 8. 1901 – 21. 8. 1901?] in Welsberg-Taisten}\toendnotes[C]{\smallbreak}
\Standort{CUL, Schnitzler, B 96a.}
\physDesc{Bildpostkarte, 428 Zeichen
\newline{}Handschrift: schwarze Tinte, deutsche Kurrent
\newline{}Versand: Stempel: »\nobreak{}\oindex{I., Innere Stadt@\textbf{I., Innere Stadt}, \emph{Verwaltungsgebiet}|pwk}Wien 1/1 1, 16. 8. 01, 8–9N\nobreak{}«.  }\toendnotes[C]{\smallbreak}\pstart{}{\pb}Herrn D\textsuperscript{r}\textsc{Arthur Schnitzler}\pend{}\pstart{}Bad \textsc{Waldbrunn} bei \textsc{Welsberg}\oindex{Wildbad Waldbrunn@\textbf{Wildbad Waldbrunn}, \emph{Spa}|pw}\pend{}\pstart{}(Puſterthal\oindex{Pustertal@\textbf{Pustertal}, \emph{Tal}|pw})\pend{}{\bigskip}
\pstart
           \noindent{}\centering{}{\pb}\textcolor{gray}{\textbf{Gruss aus Wien\oindex{Wien@\textbf{Wien}, \emph{Verwaltungsgebiet}|pw}}}\pend
           
\pstart
           \centering{}\textcolor{gray}{\textbf{Stock im Eisen\oindex{Wien@\textbf{Wien}!I., Innere Stadt@\textbf{I., Innere Stadt}!Stock im Eisen@\textbf{Stock im Eisen}, \emph{Monument}|pw}.}}\pend
           \vspace{1em}
\pstart{}{\pb}Lieber Arthur!\pend\vspace{0.5em}
\pstart
           Zur Abwechſlung{ }ſchicke einmal ich Ihnen eine Anſichtskarte aus meiner So{\geminationm}ergrüße, in der ich jetzt wol{ }ſchon bleiben. werde.
               Beſten Dank für \label{K_L04303-1v}\edtext{Brief}{\lemma{\textnormal{\emph{Brief}}}\Cendnote{\textnormal{XXXX Auszeichnungsfehler: Dokument L04078 nicht gefunden.
               }}}\label{K_L04303-1} und \label{K_L04303-2v}\edtext{Karte}{\lemma{\textnormal{\emph{Karte}}}\Cendnote{\textnormal{{XXXX ref} XXXX
               }}}\label{K_L04303-2}. Mitzuteilen habe ich rein gar nichts. Es geht
               unglaublich wenig vor. Nicht einmal ein Feuilleton von Sil
                  Vara\pwindex{Silberer, Geza 1.\,12.\,1876 Vršac – 5.\,4.\,1938 Wien@\textsc{Silberer, Geza} (1.\,12.\,1876 Vršac – 5.\,4.\,1938 Wien), \emph{Schriftsteller, Journalist}|pw}. Bitte, beſtellen Sie meine{ }ſchönſten Grüße den Damen\pwindex{Schnitzler, Olga 17.\,1.\,1882 Wien – 13.\,1.\,1970 Lugano@\textsc{Schnitzler, Olga} (17.\,1.\,1882 Wien – 13.\,1.\,1970 Lugano), \emph{Schauspielerin, Sängerin}|pwv}\pwindex{Steinrück, Elisabeth 19.\,11.\,1885 – 7.\,4.\,1920 Partenkirchen@\textsc{Steinrück, Elisabeth} (19.\,11.\,1885 – 7.\,4.\,1920 Partenkirchen)|pwv} und D\textsuperscript{r}{ }\textsc{Goldma{\geminationn}\pwindex{Goldmann, Paul 31.\,1.\,1865 Breslau – 25.\,9.\,1935 Wien@\textsc{Goldmann, Paul} (31.\,1.\,1865 Breslau – 25.\,9.\,1935 Wien), \emph{Schriftsteller, Journalist}|pw}}\pend
           
\pstart
           Herzlichſt Ihr{\\[\baselineskip]}\spacefill\mbox{Gustav}\pend
           \leftskip=0em{}\selectlanguage{ngerman}\endnumbering\briefempfaengerindex{Schnitzler, Arthur@\textsc{Schnitzler, Arthur}!zzzSchwarzkopf, Gustav@\emph{von Gustav Schwarzkopf}!1901-08-161@{16. 8. 1901}|)be}\mylabel{L04303h}
\begin{anhang}
\end{anhang}\newcommand{\dateiname}{L04303}\newcommand{\titel}{Gustav Schwarzkopf an Arthur Schnitzler, 16. 8. 1901}\newcommand{\editorInnen}{Herausgegeben von Jahnke, SelmaMüller, Martin Anton}\input{../tex-inputs/latex-leseansicht-abspann}
